\documentclass[10pt,a4paper]{amsart}
\usepackage[margin=19mm]{geometry}
\usepackage{amsmath,amssymb,mathtools}
\usepackage[scr=rsfs,cal=euler]{mathalfa}
\usepackage{xfrac}
\usepackage[unicode,hidelinks,bookmarks=false]{hyperref}
\newcommand{\ie}{i.e.\ }
\newcommand{\cf}{cf.\ }
\newcommand{\resp}{respectively\ }
\newcommand{\Subsection}{\subsection}
\providecommand{\nobreakdash}{\nobreak\hbox{-}}
\setlength{\emergencystretch}{2em}
\multlinegap0pt
\usepackage{amssymb}
\usepackage{xcolor}
\usepackage{cancel}
\usepackage[shortlabels]{enumitem}
\newtheorem{lemma}{Lemma}
\DeclareMathOperator{\Config}{Config}
\DeclareMathOperator{\Table}{Table}
\DeclareMathOperator{\coord}{coord}
\DeclareMathOperator{\flag}{flag}
\newcommand{\pd}{\mathsf{pd}}
\newcommand{\tuple}{\mathbf}
\title{Lower Bounds on Inverse Cellular Automata via~Proof Complexity}
\author{Maryia Kapytka}
\date{Source: arXiv:2604.01041v1 (CC BY 4.0)}
\begin{document}
\maketitle

\noindent\emph{Source and licence.} Lower Bounds on Inverse Cellular Automata via~Proof Complexity. Version arXiv:2604.01041v1 (CC BY 4.0), distributed by arXiv under the Creative Commons Attribution 4.0 International licence, http://creativecommons.org/licenses/by/4.0/, as recorded in the arXiv licence metadata for this exact version and shown on the abstract page https://arxiv.org/abs/2604.01041v1. Full licence notice: \texttt{licenses/arxiv-2604.01041-CC-BY-4.0.txt}.\par\smallskip

\noindent\textit{Editorial scope.} Counted unit: the article's lemma lm:new\_lemma (source0.tex lines 1355--1361). The article's complete original proof of it is delivered with it (source0.tex lines 1363--1478, one \texttt{proof} environment of 116 lines). No mathematics is written, repaired or re-derived: the statement and the proof are the article's own lines, in the article's own notation, and no retained line is substituted or re-flowed.  No further source-local context is needed: every cross-reference of every retained line resolves inside the two delivered spans.  Nothing was omitted from any retained span, so the package carries no omission record.  No retained line contains a citation, so the wrapper carries no bibliography and no bibliography entry.  No retained line contains a figure environment, an image-inclusion command or drawing code, so no image is delivered and the package declares no asset dependency.  Editorial formatting only, in addition: an AMS \texttt{amsart} preamble that restates the article's own declarations and macros used by the retained lines, namely the environments \texttt{lemma} on separate counters, together with the standard packages those lines need.  The retained environments print in this document as Lemma 1 in order of appearance, whereas the article prints the same items with the numbering of its own full document.  The complete native source of the article as distributed in the arXiv e-print for this exact version is bundled unchanged as \texttt{sources/197632/source0.tex}.
\begin{lemma}\label{lm:new_lemma}
Let $\tuple a=(a_1,a_2,\dots,a_m)\in\{0,1\}^m$ and let
$C\in\Config_{n,m}$.
If $C$ is locally correct and, for every $j\in\{1,2,\dots,m\}$, every cell
with state $s$ satisfying $\coord(s)=(0,j)$ also satisfies $a(s)=a_j$, then
$C$ is similar to $\Table_\varphi(\tuple a)$.
\end{lemma}

\begin{proof}
We show that every component of the state of each cell apart from label coincides with the corresponding entry
of $\Table_\varphi(\tuple a)$.

\paragraph{Coordinates.}
We show that for every cell $c=(i,j,s)$ with 
$(i,j)\in [n+1]\times[m+2]$ it holds
\[
\coord(s)=(i,j).
\]

We first prove correctness of the row index by induction on $i$,
using local correctness rule~(A).

\medskip
\noindent
\emph{Base case ($i=0$).}
Let $c=(0,j,s)$ for some $j\in[m+2]$, and let $c'$ be the cell above
$c$. Since $C\in\Config_{n,m}$, the cell $c'$ is outside the finite
part of the table and therefore is in the quiescent state $\tuple q$.
By index consistency~(A), a cell whose above neighbour is in state
$\tuple q$ must have first coordinate $0$. Hence
\[
\coord(s)=(0,j)
\]
for some $j$, and therefore the row index is correct.

\medskip
\noindent
\emph{Induction step.}
Assume that for some $i\ge 0$ every cell in rows $0, 1, \dots, i$ has correct
first coordinate in its state. 

Let $c=(i+1,j,s)$ and let $c'=(i,j,s')$ be the
cell above $c$.  By inductive assumption $\coord(s') = (i,j)$. By rule~(A), as $c$ is below neighbour of $c'$, $\coord(s) =(i+1, j)$. 
Thus every cell in row $i+1$ has the correct first coordinate.
This completes the induction on $i$.

An analogous induction on the column index $j$ shows that the second
coordinate is also correct. 

In the remainder of the proof, we rely on each cell’s state containing its correct coordinates.


\paragraph{Assignment values.}
By assumption, the $0$th row contains the correct assignment values $a_j$.
By vertical consistency~(C1), these values are propagated unchanged to all cells
below the $0$th row. By an easy induction on row index, every cell with coordinates $(i,j)$ for $i \in \{1, 2, \dots, n\}$ and $j \in \{1, 2, \dots, m\}$ contains the
correct value $a_j$.

\paragraph{Partial disjunctions.}
We show that for every $i\in\{1,\dots,n\}$ and every $j\in\{1,\dots,m\}$,
every cell with state $s$ such that $\coord(s)=(i,j)$ satisfies
\[
\mathsf{pd}(s)=\bigl(\bigvee_{u\le j} C_{i,u}\bigr)(\tuple a).
\]
The proof is by induction on~$j$.

\emph{Base case $j=1$.}
By rule~(D2), $\mathsf{pd}(s)=1$ if and only if either
$\flag(s)=1$ and $a_1=1$, or $\flag(s)=-1$ and $a_1=0$.
In both cases, the literal in column~$1$ satisfies clause~$C_i$,
so $C_{i,1}(\tuple a)=1$.
If none of these cases holds, then either $x_1$ does not satisfy $C_i$
or does not occur in it, and hence $C_{i,1}(\tuple a)=0$.
Thus $\mathsf{pd}(s)=C_{i,1}(\tuple a)$.

\emph{Induction step.}
Assume the claim holds for column $k-1$, and let $s$ be a state such that
$\coord(s) = (i,k)$.  
Let $s'$ be the state of its left neighbour.

Assume first $\mathsf{pd}(s')=1$. (D1) implies that $\mathsf{pd}(s)=1$. Moreover, by inductive assumption $\bigl(\bigvee_{u \leq k-1} C_{i,u}\bigr)(\tuple a) = \pd(s') = 1$, which means  
\[
\bigl(\bigvee_{u \leq k} C_{i,u}\bigr)(\tuple a) = \bigl(\bigvee_{u \leq k-1} C_{i,u}\bigr)(\tuple a) \vee C_{i, k} (\tuple a)= 1 \vee C_{i, k} (\tuple a) =1.
\]
 Therefore $\mathsf{pd}(s) = \bigl(\bigvee_{u \leq k} C_{i,u}\bigr)(\tuple a)$ in this case.

Assume $\pd(s')=0$. Here (D1) yields that $\mathsf{pd}(s)=1$ if and only if $\flag(s)=1\wedge a(s)=1$ or $\flag(s)=-1\wedge a(s)=0$. Similarly as in the base step, if $\mathsf{pd}(s)=1$, then in both cases $a_{k}$ satisfies $C_i$ and 
 \[
 \bigl(\bigvee_{u \leq k} C_{i,u}\bigr)(\tuple a) = \bigl(\bigvee_{u \leq k-1} C_{i,u}\bigr)(\tuple a) \vee C_{i,k}(\tuple a) = \bigl(\bigvee_{u \leq k-1} C_{i,u}\bigr)(\tuple a) \vee 1=1.
 \]
   If $\mathsf{pd}(s)=0$, then either $a_{k}$ does not satisfy or does not belong to $C_i$. In both of these cases, as $\bigl(\bigvee_{u \leq k-1} C_{i,u}\bigr)(\tuple a) =0$ by inductive assumption, then 
    \[
 \bigl(\bigvee_{u \leq k} C_{i,u}\bigr)(\tuple a) = \bigl(\bigvee_{u \leq k-1} C_{i,u}\bigr)(\tuple a) \vee C_{i,k}(\tuple a) = 0 \vee 0=0,
 \]
   
 and we are done.

\paragraph{Partial conjunctions.}
Finally, we show that for every $i\in\{1,\dots,n\}$ and every cell with $s$ and
$\coord(s)=(i,m+1)$,
\[
\mathsf{pc}(s)=\bigl(\bigwedge_{v\le i} C_v\bigr)(\tuple a).
\]
The proof is by induction on~$i$.

\emph{Base case $i=1$.}
By (C3) $\mathsf{pc}(s)=1$ if and only if $\mathsf{pd}(s')=1$ for the left neighbour with state $s'$. We have already shown that $\mathsf{pd}(s')=1$ if and only if $\bigl(\bigvee_{u \leq m} C_{1,u}\bigr)(\tuple a) = 1$. It remains to notice that $(\bigvee_{u \leq m} C_{1,u}\bigr)$ is exactly $C_1$ and we are done.


\emph{Induction step.}
Assume the claim holds for row $l-1$, and let $s$ be a state such that
$\coord(s)=(l,m+1)$. Let $s''$ be the state of the above neighbour.

If $\mathsf{pc}(s'')=0$, then according to (C2), $\mathsf{pc}(s)=0$. By inductive assumption $\bigl(\bigwedge_{v \leq l-1}C_v\bigr) (\tuple a) =\mathsf{pc}(s)=0$, therefore
\[
\bigl(\bigwedge_{v \leq l}C_v\bigr) (\tuple a) = \bigl(\bigwedge_{v \leq l-1}C_v\bigr) (\tuple a) \wedge C_{l}(\tuple a) = 0 \wedge C_{l}(\tuple a) =0.
\]

Assume $\mathsf{pc}(s'')=1$. According to (C2) $\mathsf{pc}(s)=1$ if and only if for left neighbour of $c$ with state $s'$ it holds $\mathsf{pd}(s')=1$. Similarly as in the base case, we have shown in the first part of the proof that $\mathsf{pd}(s') = (\bigvee_{u \leq m} C_{l,u}\bigr)(\tuple a) = C_{l}(\tuple a)$. So $\mathsf{pc}(s)=1$ if and only if $C_{l}(\tuple a) =1$. Then 
\[
\bigl(\bigwedge_{v \leq l}C_v\bigr) (\tuple a) = \bigl(\bigwedge_{v \leq l-1}C_v\bigr) (\tuple a) \wedge C_{l}(\tuple a) = 1 \wedge  C_{l}(\tuple a) = C_{l}(\tuple a).
\]
This shows that $\mathsf{pc}(s)=1$ if and only if $\bigl(\bigwedge_{v \leq l}C_v\bigr) (\tuple a) =1$ and we are done.
\end{proof}

\end{document}
